%% kapa: exam=1 round=19 year=2008 period=1 form=A subject=civil_law subject_ko=민법 seq=37 number=37
%% source: Q-Net 감정평가사 기출문제 (cst00309) artlSeq=1001962 file=extracted/1008/1001962/2008년 제19회 감정평가사 1차시험 기출문제/1교시.hwp
%% conversion: kordoc 4.24.51 -> tools/gen1.py (auto); review=proofread-sonnet; flags=none
\begin{question}{37}
저당토지 위의 건물에 대한 일괄경매청구권에 관한 설명으로 옳은 것은? (다툼이 있는 경우에는 판례에 의함)
\begin{choices}
  \item 일괄경매청구권은 원칙적으로 저당권설정자가 건물을 축조하여 소유하고 있는 경우에 한한다.
  \item 저당토지 위의 건물에 대한 일괄경매청구권의 요건을 구비한 경우 저당권자는 일괄경매를 신청할 의무를 부담한다.
  \item 일괄경매의 대상인 건물은 원칙적으로 저당토지 위에 신축되어야 하므로 증축한 건물의 일부가 다른 토지에 걸쳐 있는 경우라면 건물 전부에 대하여 일괄경매를 청구할 수 없다.
  \item 토지 저당권자는 건물의 매각대금에 대해서도 우선변제를 받을 수 있다.
  \item 저당토지만 경매해도 충분히 피담보채권을 변제 받을 수 있는 경우라면 건물에 대한 일괄경매가 허용되지 않는다.
\end{choices}
\end{question}
\answer{1}
